\honest{Fake Optimization Criteria}{Eliezer Yudkowsky}{2007}

I have written at length, with a link on nearly every word, about how our beliefs seem to fit the evidence better than they do. If we could beat this, I say, ``we would be able to recover from almost any other error of fact.'' I do not argue for this.

The same problem arises for decisions: seeing what a criterion really endorses. We seem able to argue that any goal implies any action. ``There's no mystery about this'': we are descended from the ancestors who argued best that the good of the tribe required killing their rival Uglak. \nb{The talent was hedged with ``seem''; its evolutionary explanation is stated as settled, with no evidence.}

Can we prove that Robert Mugabe, if he cared only for Zimbabwe, would resign? Can we even prove that people rationalize, when we never know exactly what a criterion implies? In a parenthesis I answer the second question myself: people's views on office politics track their interests. What I really want is an optimizer with a known goal and no other desires, to compare with human reasoning about the same goal.

Natural selection is that optimizer. It has no mercy and no politics, and its output is optimized ``only for inclusive genetic fitness.'' So we can compare. What did Wynne-Edwards expect group selection for small populations to produce? ``Voluntary individual restraint in breeding, and enough food for everyone.'' What did the laboratory produce? ``Cannibalism.'' \nb{The editor found no such prediction by Wynne-Edwards, whose book came fourteen years before the experiment. See ``The Tragedy of Group Selectionism.''}

Does human morality care about fitness? No. But if we cannot hear one note clearly, I ask, how will we hear an orchestra? We need simple practice cases. \nb{The post gives no evidence that practice on evolutionary cases helps anyone judge moral principles. In its own example the answer came from an experiment, and for human criteria there is no experiment.}
